\section{2026 Geometry \& Topology}

\subsection{Exam}

\begin{exercise}
\label{geometry:2026:1}
Let $S \subset \mathbb{R}^3$ be a smooth regular surface without boundary, and let
$p \in S$. Let $\{A_{\epsilon}\}_{\epsilon>0}$ be a family of regions in $S$ such
that each $A_{\epsilon}$ contains $p$, has area $|A_{\epsilon}|$, and shrinks to
$\{p\}$ as $\epsilon \to 0$. Let $N:S\to S^2$ be the Gauss map. Show that the
Gaussian curvature at $p$ satisfies
\[
    |K(p)|=\lim_{\epsilon\to 0}
    \frac{\operatorname{Area}(N(A_{\epsilon}))}
         {\operatorname{Area}(A_{\epsilon})}.
\]
\end{exercise}

\begin{solution}
The differential of the Gauss map is the negative of the shape operator:
\[
    dN_q=-S_q:T_qS\longrightarrow T_{N(q)}S^2.
\]
Therefore its Jacobian determinant is
\[
    |\det dN_q|=|\det S_q|=|K(q)|.
\]

First assume $K(p)\neq 0$. Then $dN_p$ is nonsingular, so by the inverse function
theorem $N$ is a diffeomorphism from a sufficiently small neighborhood $U$ of $p$
onto its image. For all sufficiently small $\epsilon$ we have $A_{\epsilon}\subset U$.
The change-of-variables formula gives
\[
    \operatorname{Area}(N(A_{\epsilon}))
    =\int_{A_{\epsilon}} |\det dN_q|\,dA(q)
    =\int_{A_{\epsilon}} |K(q)|\,dA(q).
\]
Thus
\[
    \frac{\operatorname{Area}(N(A_{\epsilon}))}
         {\operatorname{Area}(A_{\epsilon})}
    =
    \frac{1}{\operatorname{Area}(A_{\epsilon})}
    \int_{A_{\epsilon}} |K(q)|\,dA(q).
\]
Since $K$ is continuous and the regions $A_{\epsilon}$ shrink to $p$, the average
of $|K|$ over $A_{\epsilon}$ tends to $|K(p)|$.

Now assume $K(p)=0$. The area formula gives
\[
    \operatorname{Area}(N(A_{\epsilon}))
    \leq \int_{A_{\epsilon}} |K(q)|\,dA(q),
\]
where possible multiplicity of the Gauss map can only increase the integral on the
right. Hence
\[
    0\leq
    \frac{\operatorname{Area}(N(A_{\epsilon}))}
         {\operatorname{Area}(A_{\epsilon})}
    \leq
    \frac{1}{\operatorname{Area}(A_{\epsilon})}
    \int_{A_{\epsilon}} |K(q)|\,dA(q)\longrightarrow |K(p)|=0.
\]
This proves the desired formula in all cases.
\end{solution}

\begin{exercise}
\label{geometry:2026:2}
Let $(M^2,g)$ be a smooth, compact, oriented two-dimensional Riemannian manifold
with nonempty boundary $\partial M$. Assume that
\begin{enumerate}
    \item the Gaussian curvature satisfies $K_g\geq 0$ on $M$, and
    \item the geodesic curvature of the boundary satisfies $k_g\geq 1$ along
    $\partial M$. Here $k_g$ denotes the geodesic curvature of $\partial M$,
    computed with respect to the outward pointing unit normal along the boundary.
\end{enumerate}
Suppose moreover that
\[
    \operatorname{Length}(\partial M)\geq 2\pi.
\]
Show that $(M,g)$ is isometric to the unit disk $(D^2(1),g_{\mathrm{flat}})$
equipped with the standard Euclidean metric.
\end{exercise}

\begin{solution}
We use the boundary convention for which the Euclidean unit disk has outward
geodesic curvature $k_g=1$. With this convention Gauss--Bonnet reads
\[
    \int_M K_g\,dA+\int_{\partial M} k_g\,ds=2\pi\chi(M).
\]
By the hypotheses,
\[
    2\pi\chi(M)
    =
    \int_M K_g\,dA+\int_{\partial M} k_g\,ds
    \geq
    \operatorname{Length}(\partial M)
    \geq 2\pi.
\]
Hence $\chi(M)\geq 1$.

A connected compact oriented surface with genus $g$ and $b\geq 1$ boundary
components has
\[
    \chi(M)=2-2g-b\leq 1,
\]
with equality only when $g=0$ and $b=1$. Therefore $M$ is topologically a disk and
$\chi(M)=1$. Since equality must hold in the chain of inequalities above, we get
\[
    \int_M K_g\,dA=0,\qquad
    \int_{\partial M}(k_g-1)\,ds=0,\qquad
    \operatorname{Length}(\partial M)=2\pi.
\]
The functions $K_g$ and $k_g-1$ are continuous and nonnegative, so
\[
    K_g\equiv 0\quad\text{on }M,\qquad
    k_g\equiv 1\quad\text{on }\partial M.
\]

It remains to identify the flat disk. Since $M$ is simply connected and flat, the
Levi-Civita connection has trivial holonomy. Choose a parallel orthonormal frame
and integrate its dual coframe; this gives a developing map
\[
    F:(M,g)\longrightarrow \mathbb{R}^2
\]
which is a local isometry. The boundary curve $F(\partial M)$ has length $2\pi$
and constant Euclidean curvature $1$ with respect to the corresponding outward
normal convention. A closed plane curve parameterized by arclength with constant
curvature $1$ is a unit circle, traversed once.

Thus $F$ maps $\partial M$ onto the boundary of a Euclidean unit disk. Since $F$ is
a local isometry on the interior and has degree one relative to the boundary circle,
it maps $M$ isometrically onto the disk bounded by that circle. Therefore
$(M,g)$ is isometric to $(D^2(1),g_{\mathrm{flat}})$.
\end{solution}

\begin{exercise}
\label{geometry:2026:3}
Let $G$ be a connected $n$-dimensional Lie group equipped with a bi-invariant
Riemannian metric $\langle\cdot,\cdot\rangle$.
\begin{enumerate}
    \item Show that the sectional curvature of $G$ is nonnegative.
    \item Assume that the Lie algebra $\mathfrak{g}$ of $G$ has trivial center, i.e.,
    \[
        \mathfrak{z}(\mathfrak{g})
        :=\{X\in\mathfrak{g}\mid [X,Y]=0\text{ for all }Y\in\mathfrak{g}\}
        =\{0\}.
    \]
    Show that $G$ is compact.
    \item Suppose that $G$ is simply connected. Show that $G$ decomposes as a direct
    product
    \[
        G\cong G_0\times \mathbb{R}^k,
    \]
    where $G_0$ is a simply connected compact Lie group whose Lie algebra has
    trivial center, and $\mathbb{R}^k$ is the additive Lie group.
\end{enumerate}
\end{exercise}

\begin{solution}
For a bi-invariant metric, the inner product on $\mathfrak{g}$ is
$\operatorname{Ad}$-invariant. Equivalently,
\[
    \langle [X,Y],Z\rangle+\langle Y,[X,Z]\rangle=0
    \qquad
    (X,Y,Z\in\mathfrak{g}).
\]
For left-invariant vector fields the Koszul formula gives
\[
    \nabla_XY=\frac{1}{2}[X,Y].
\]
With the curvature sign convention
$R(X,Y)Z=\nabla_X\nabla_YZ-\nabla_Y\nabla_XZ-\nabla_{[X,Y]}Z$, this gives
\[
    R(X,Y)Z=-\frac{1}{4}[[X,Y],Z],
\]
and $K(X,Y)=\langle R(X,Y)Y,X\rangle$ for an orthonormal pair. Using invariance of
the inner product,
\[
    \langle R(X,Y)Y,X\rangle
    =
    \frac{1}{4}\|[X,Y]\|^2.
\]
Thus every sectional curvature is nonnegative:
\[
    K(X,Y)=\frac{1}{4}\|[X,Y]\|^2\geq 0.
\]

Now assume $\mathfrak{z}(\mathfrak{g})=0$. Since every $\operatorname{ad}_X$ is
skew-adjoint, the Killing form satisfies
\[
    B(X,X)=\operatorname{tr}(\operatorname{ad}_X^2)\leq 0.
\]
Its radical is the center for a Lie algebra carrying a positive definite
ad-invariant inner product. Since the center is zero, $-B$ is positive definite;
therefore $\mathfrak{g}$ is a compact semisimple Lie algebra. The simply connected
Lie group with this Lie algebra is compact, and every connected Lie group with
this Lie algebra is a quotient of it by a finite central subgroup. Hence $G$ is
compact.

Finally suppose $G$ is simply connected. Let
$\mathfrak{z}=\mathfrak{z}(\mathfrak{g})$. Because the metric is ad-invariant, the
orthogonal complement $\mathfrak{z}^{\perp}$ is an ideal: if $X\in\mathfrak{z}$,
$Y\in\mathfrak{z}^{\perp}$, and $Z\in\mathfrak{g}$, then
\[
    \langle X,[Z,Y]\rangle=\langle [X,Z],Y\rangle=0.
\]
Moreover,
\[
    \mathfrak{g}=\mathfrak{z}\oplus \mathfrak{z}^{\perp}
\]
as a direct sum of commuting ideals. The ideal $\mathfrak{z}$ is abelian, so its
simply connected Lie group is $\mathbb{R}^k$, where $k=\dim\mathfrak{z}$. The ideal
$\mathfrak{z}^{\perp}$ has trivial center; by the previous part its simply
connected Lie group $G_0$ is compact. Since $G$ is simply connected and the two
ideals commute, the corresponding connected subgroups multiply to a global direct
product:
\[
    G\cong G_0\times \mathbb{R}^k.
\]
\end{solution}

\begin{exercise}
\label{geometry:2026:4}
Let $S^{2026}$ be the unit sphere in $\mathbb{R}^{2027}$.
\begin{enumerate}
    \item Compute the Euler class of the tangent bundle
    $TS^{2026}\to S^{2026}$ with $\mathbb{Z}_2$ coefficients.
    \item Let $US^{2026}\to S^{2026}$ denote the unit tangent sphere bundle; that is,
    the fiber over each point is identified with $S^{2025}$. Compute the Poincare
    series
    \[
        P_t(US^{2026};\mathbb{Z}_2)
        =
        \sum_{i\geq 0} b_i(US^{2026};\mathbb{Z}_2)t^i,
    \]
    where $b_i(US^{2026};\mathbb{Z}_2)$ denotes the $i$-th Betti number.
    \item Let $p,q\in S^{2026}$ be two non-antipodal points. Consider the space of
    piecewise smooth paths from $p$ to $q$,
    \[
        \Omega(S^{2026};p,q)
        =
        \{\gamma:[0,1]\to S^{2026}\mid \gamma(0)=p,\ \gamma(1)=q\},
    \]
    equipped with the energy functional
    \[
        E(\gamma)=\frac{1}{2}\int_0^1 |\dot{\gamma}(t)|^2\,dt.
    \]
    \begin{enumerate}
        \item Determine all critical points of $E$.
        \item The index $\lambda(E,\gamma)$ of $E$ at a critical point $\gamma$ is
        defined as the dimension of the maximal subspace of
        $T_{\gamma}\Omega(S^{2026};p,q)$ on which $\operatorname{Hess}E$ is negative
        definite. Compute $\lambda(E,\gamma)$.
        \item Determine the homotopy type of the based loop space
        $\Omega(S^{2026},p)$.
        \item Compute the Poincare series of both the based loop space
        $\Omega(S^{2026},p)$ and the free loop space $\Lambda S^{2026}$ with
        $\mathbb{Z}_2$ coefficients.
    \end{enumerate}
\end{enumerate}
\end{exercise}

\begin{solution}
Let $u$ denote the generator of
$H^{2026}(S^{2026};\mathbb{Z})\cong\mathbb{Z}$. The integral Euler class of the
tangent bundle satisfies
\[
    \langle e(TS^{2026}),[S^{2026}]\rangle
    =
    \chi(S^{2026})=2,
\]
so $e(TS^{2026})=2u$. Reducing modulo $2$ gives
\[
    e(TS^{2026};\mathbb{Z}_2)=0
    \in H^{2026}(S^{2026};\mathbb{Z}_2).
\]
Equivalently, this is the top Stiefel--Whitney class
$w_{2026}(TS^{2026})$, and it vanishes because the Euler characteristic is even.

For the unit tangent sphere bundle
\[
    S^{2025}\longrightarrow US^{2026}\longrightarrow S^{2026},
\]
the Serre spectral sequence with $\mathbb{Z}_2$ coefficients has only four
potentially nonzero entries:
\[
    E_2^{0,0},\quad E_2^{2026,0},\quad E_2^{0,2025},\quad E_2^{2026,2025}.
\]
The only possible nonzero differential is the transgression
\[
    d_{2026}:E_{2026}^{0,2025}\longrightarrow E_{2026}^{2026,0},
\]
and this differential is multiplication by the mod-$2$ Euler class of
$TS^{2026}$. Since that Euler class is zero, the spectral sequence collapses.
Therefore the nonzero Betti numbers occur in degrees
\[
    0,\quad 2025,\quad 2026,\quad 4051,
\]
each with value $1$. Hence
\[
    P_t(US^{2026};\mathbb{Z}_2)
    =
    1+t^{2025}+t^{2026}+t^{4051}.
\]

Now consider the energy functional. Put $N=2026$. First assume $p\neq q$, and let
\[
    \alpha=d(p,q)\in(0,\pi).
\]
Choose the unit vector $v\in T_pS^N$ pointing from $p$ to $q$ along the minimizing
geodesic, so
\[
    q=\cos\alpha\,p+\sin\alpha\,v.
\]
Let
\[
    c(s)=\cos s\,p+\sin s\,v
\]
be the corresponding unit-speed great circle. The critical points of $E$ are
exactly the constant-speed geodesics from $p$ to $q$. Since $p$ and $q$ are not
antipodal, all such geodesics lie on this great circle and are
\[
    \gamma_m^+(t)=c((\alpha+2\pi m)t),
    \qquad m=0,1,2,\ldots,
\]
and
\[
    \gamma_m^-(t)=c(-(2\pi-\alpha+2\pi m)t),
    \qquad m=0,1,2,\ldots.
\]
The lengths are
\[
    L_m^+=\alpha+2\pi m,
    \qquad
    L_m^-=2\pi-\alpha+2\pi m.
\]

By the Morse index theorem for geodesics with fixed endpoints, the index is the
sum of the multiplicities of conjugate points to $p$ along the geodesic before the
endpoint. On the unit sphere $S^N$, conjugate points along any geodesic occur at
distances
\[
    \pi,2\pi,3\pi,\ldots
\]
and each has multiplicity $N-1=2025$. Since $\alpha\in(0,\pi)$, the endpoint is not
conjugate to $p$. Therefore
\[
    \lambda(E,\gamma_m^+)
    =
    2025\cdot \#\{j\geq 1:j\pi<\alpha+2\pi m\}
    =
    2025\cdot 2m,
\]
and
\[
    \lambda(E,\gamma_m^-)
    =
    2025\cdot \#\{j\geq 1:j\pi<2\pi-\alpha+2\pi m\}
    =
    2025(2m+1).
\]
If $p=q$, the same endpoint problem becomes Morse--Bott rather than Morse. The
critical points are the constant loop, with index $0$, and for each $m\geq 1$ the
closed geodesics
\[
    \gamma_{m,v}(t)=\cos(2\pi mt)p+\sin(2\pi mt)v,
    \qquad v\in U_pS^N.
\]
These form a critical manifold diffeomorphic to $S^{N-1}$. The endpoint is
conjugate, so the Hessian has a kernel tangent to this critical manifold, and the
negative index is
\[
    \lambda(E,\gamma_{m,v})=(N-1)(2m-1)=2025(2m-1).
\]

The based loop space of a sphere is described by the James construction:
\[
    \Omega S^{2026}\simeq J(S^{2025}).
\]
The precise theorem being used is the James weak equivalence
$J(X)\simeq \Omega\Sigma X$ for connected CW complexes $X$. Since
$S^{2026}=\Sigma S^{2025}$, this applies with $X=S^{2025}$. In particular
$J(S^{2025})$ has one cell for each word length in the generator of
$S^{2025}$, so it has the homotopy type of a CW complex with one cell in each
dimension
\[
    0,\ 2025,\ 2\cdot 2025,\ 3\cdot 2025,\ldots.
\]
Equivalently, with $\mathbb{Z}_2$ coefficients,
$H_*(\Omega S^{2026};\mathbb{Z}_2)$ is the tensor algebra on one generator of
degree $2025$. Hence
\[
    P_t(\Omega(S^{2026},p);\mathbb{Z}_2)
    =
    \frac{1}{1-t^{2025}}.
\]

For the free loop space, use the evaluation fibration
\[
    \Omega S^{2026}\longrightarrow \Lambda S^{2026}\longrightarrow S^{2026}.
\]
It has a section given by constant loops. In the Serre spectral sequence over
$\mathbb{Z}_2$, the base has only the rows in degrees $0$ and $2026$, and the
section forces the base fundamental class to survive. Multiplicativity then
forces the only possible transgression to vanish on the algebra generated by the
loop-space class. Thus the spectral sequence collapses as a vector-space
calculation, and
\[
    H_*(\Lambda S^{2026};\mathbb{Z}_2)
    \cong
    H_*(S^{2026};\mathbb{Z}_2)\otimes
    H_*(\Omega S^{2026};\mathbb{Z}_2)
\]
as graded $\mathbb{Z}_2$-vector spaces, and
\[
    P_t(\Lambda S^{2026};\mathbb{Z}_2)
    =
    \frac{1+t^{2026}}{1-t^{2025}}.
\]
\end{solution}

\begin{exercise}
\label{geometry:2026:5}
Let $M^n$ be a connected, closed, smooth, aspherical manifold of dimension
$n\geq 1$, that is, its universal covering space $\widetilde{M}$ is contractible.
\begin{enumerate}
    \item Prove that the universal cover $\widetilde{M}$ is noncompact.
    \item Prove that every nontrivial element of $\pi_1(M)$ has infinite order.
\end{enumerate}
\end{exercise}

\begin{solution}
Suppose first that $\widetilde{M}$ were compact. Since the fiber of the covering
map $\widetilde{M}\to M$ is a discrete closed subset of the compact space
$\widetilde{M}$, each fiber would be finite. Thus $\widetilde{M}\to M$ would be a
finite-sheeted covering. In particular $\widetilde{M}$ would be a closed
$n$-manifold.

But a connected closed $n$-manifold has
\[
    H_n(\widetilde{M};\mathbb{Z}_2)\cong \mathbb{Z}_2,
\]
whereas a contractible space has zero homology in every positive degree. Since
$n\geq 1$, this is impossible. Therefore $\widetilde{M}$ is noncompact.

For the second part, the group $\pi_1(M)$ acts on $\widetilde{M}$ by deck
transformations. This action is free: a deck transformation with a fixed point is
the identity.

Assume that some nontrivial element $g\in\pi_1(M)$ has finite order. Then some
power of $g$ has prime order $p$; call that element $h$. The cyclic group
$\langle h\rangle\cong\mathbb{Z}/p\mathbb{Z}$ acts smoothly on the contractible
manifold $\widetilde{M}$. Since $\widetilde{M}$ is contractible, it is
$\mathbb{Z}_p$-acyclic. The Smith fixed point theorem applies in the following
form: a smooth action of a finite $p$-group on a finite-dimensional
$\mathbb{Z}_p$-acyclic manifold has fixed point set again
$\mathbb{Z}_p$-acyclic, in particular nonempty. Hence $h$ must fix some point of
$\widetilde{M}$.

This contradicts freeness of the deck transformation action. Hence no nontrivial
element of $\pi_1(M)$ can have finite order.
\end{solution}

\begin{exercise}
\label{geometry:2026:6}
Let $(M^n,g)$ be a complete Riemannian manifold of dimension $n\geq 2$. Suppose
that there exists $\epsilon>\frac{1}{4}$ such that
\[
    \operatorname{Ric}_g(x)\geq
    \epsilon(n-1)r(p,x)^{-2}
\]
for all $x$ with $r(p,x)=d(p,x)$ sufficiently large.
\begin{enumerate}
    \item Show that $M$ must be compact.
    \item Show, by example, that the conclusion in (1) fails if
    $\epsilon\leq \frac{1}{4}$.
\end{enumerate}
\end{exercise}

\begin{solution}
Assume for contradiction that $M$ is noncompact. Since $M$ is complete, there is a
ray
\[
    \gamma:[0,\infty)\to M,\qquad \gamma(0)=p,
\]
parameterized by arclength. Thus each segment $\gamma|_{[0,T]}$ is minimizing.
Choose $R_0$ so that
\[
    \operatorname{Ric}_g(\gamma'(t),\gamma'(t))
    \geq
    \epsilon(n-1)t^{-2}
    \qquad (t\geq R_0).
\]

Let $E_1,\ldots,E_{n-1}$ be parallel orthonormal vector fields along $\gamma$,
perpendicular to $\gamma'$. For any smooth function $\phi$ compactly supported in
$(R_0,T)$, the fields $V_i=\phi E_i$ vanish at the endpoints of
$\gamma|_{[0,T]}$. Since this segment is minimizing, the index form is
nonnegative. Summing over $i$ gives
\[
    0\leq
    \sum_{i=1}^{n-1} I(V_i,V_i)
    =
    \int_0^T
    \left((n-1)(\phi')^2
    -\phi^2\operatorname{Ric}_g(\gamma',\gamma')\right)\,dt.
\]
Using the Ricci lower bound on the support of $\phi$,
\[
    0\leq
    (n-1)\int_{R_0}^T
    \left((\phi')^2-\epsilon\frac{\phi^2}{t^2}\right)\,dt.
\]

We now choose a test function that violates this inequality. Write $t=e^s$ and set
\[
    \phi(t)=t^{1/2}\psi(\log t),
\]
where $\psi$ is a smooth compactly supported function. A direct computation gives
\[
    \int
    \left((\phi')^2-\epsilon\frac{\phi^2}{t^2}\right)\,dt
    =
    \int
    \left((\psi')^2-\left(\epsilon-\frac{1}{4}\right)\psi^2\right)\,ds,
\]
because
\[
    \phi'(t)=t^{-1/2}\left(\psi'(\log t)+\frac12\psi(\log t)\right),
\]
and the cross term $\int \psi\psi'\,ds$ vanishes for compactly supported
$\psi$. This is exactly the one-dimensional Hardy constant $1/4$.

Since $\epsilon>1/4$, choose $\psi=\psi_L$ supported in an interval of length
about $L$, equal to $1$ on the middle portion, and with endpoint cutoffs of fixed
shape. Then $\int(\psi_L')^2\,ds$ stays bounded independently of $L$, while
\[
    \left(\epsilon-\frac14\right)\int\psi_L^2\,ds
\]
grows linearly in $L$. For $L$ large the integral is negative. Translating the
support far enough to the right puts the corresponding $\phi$ inside
$(R_0,T)$ for a sufficiently large minimizing segment of the ray. This
contradicts the nonnegativity of the index form. Therefore $M$ must be compact.

The constant $1/4$ is sharp. Fix any $0<\epsilon\leq 1/4$ and choose
$\alpha\in(0,1)$ such that
\[
    \alpha(1-\alpha)=\epsilon.
\]
On $\mathbb{R}^n$ take a rotationally symmetric complete metric
\[
    g=dr^2+f(r)^2g_{S^{n-1}},
\]
where $f(r)=r$ near $0$ and $f(r)=r^{\alpha}$ for all sufficiently large $r$.
For large $r$, the radial Ricci curvature is
\[
    \operatorname{Ric}(\partial_r,\partial_r)
    =
    -(n-1)\frac{f''}{f}
    =
    (n-1)\alpha(1-\alpha)r^{-2}
    =
    \epsilon(n-1)r^{-2}.
\]
For a unit tangent vector $X$ tangent to the sphere factor, the warped product
curvature formula gives
\[
    \operatorname{Ric}(X,X)
    =
    \frac{(n-2)(1-(f')^2)-ff''}{f^2}.
\]
When $f(r)=r^{\alpha}$ this becomes
\[
    \operatorname{Ric}(X,X)
    =
    (n-2)r^{-2\alpha}
    +
    \bigl(\alpha-(n-1)\alpha^2\bigr)r^{-2}.
\]
For $n=2$ this is exactly $\alpha(1-\alpha)r^{-2}$; for $n\geq 3$ it is at least
$(n-1)\alpha(1-\alpha)r^{-2}$ for all sufficiently large $r$, because
$r^{-2\alpha}$ decays more slowly than $r^{-2}$. Hence, outside a compact set,
\[
    \operatorname{Ric}_g\geq \epsilon(n-1)r^{-2}.
\]
The metric is complete and noncompact, so compactness fails for every
$0<\epsilon\leq 1/4$. For $\epsilon\leq 0$, ordinary Euclidean space already gives
a noncompact example.
\end{solution}
